Queda por demostrar que \emph{c)} implica \emph{a)}. Razonemos por inducción
sobre $n=\dim(A)$; el teorema es trivial cuando $n=0$. Además,
es posible limitarse al caso en que $A=B$ es \emph{íntegro}. Procedamos en
dos etapas, siendo $n$ $\geqslant1$.

\medskip\noindent
I) \emph{Supongamos primero que $A$ satisface $(\mathrm S_2)$.} ---
Pongamos $X=\operatorname{Spec}(A)$, $A'=\widehat A$ y
$X'=\operatorname{Spec}(A')$, y sea $u:X'\to X$ el morfismo
canónico. Hay que demostrar que, para todo
$x'\in\operatorname{Ass}(\mathscr O_{X'})$, se tiene
$\dim(\overline{\{x'\}})=n$. Sea $f\neq0$ un elemento del ideal
maximal de $A$ y pongamos $C=A/fA$. Se sabe por (5.7.6) que
$A/fA$ satisface $(\mathrm S_1)$. Además, los ideales primos de
$A$ minimales entre los que contienen $f$ tienen altura 1 por
el Hauptidealsatz (teorema del ideal principal); como $C$ es
catenario, $C=A/fA$ es \emph{estrictamente equidimensional} y
$\dim(C)=n\setmin 1$. Se tiene
$\widehat C=C\otimes_AA'=A'/fA'$, y $f$ es $A'$-regular por
planitud (\textbf 0$_{\mathrm I}$, 6.3.4). Si se pone
$Y'=V(fA')=\operatorname{Spec}(\widehat C)$, entonces, para todo
punto maximal $y'$ de $Y'\cap\overline{\{x'\}}$, se tiene
$y'\in\operatorname{Ass}(\mathscr O_{Y'})$ por (3.4.3).
Por otra parte, $y'\neq x'$ porque $u(x')$ es el punto genérico
de $X$ (3.3.2) mientras que $u(y')\in V(fA)$. De (5.1.8) se
deduce
\begin{equation}
\dim(\overline{\{y'\}})=\dim(\overline{\{x'\}})\setmin 1.
\tag{7.2.5.1}
\label{IV.2.eq.7.2.5.1-fr}
\end{equation}

Pero el anillo cociente $C=A/fA$ satisface también la hipótesis
\emph{c)} del enunciado (5.6.1 y 5.11.7, (ii)). Por la hipótesis de
inducción, $\dim(\overline{\{y'\}})=\dim(C)=n\setmin 1$; en este caso,
la conclusión $\dim(\overline{\{x'\}})=n$ resulta de (7.2.5.1).
